% GENERATED FILE. Edit canonical lessons, then run npm run generate:books.
% canonical-source-hash: fnv1a64:64acc51645f0ea7c
% canonical-lessons: BN-C471-osukh, BN-C471-lokkhon, BN-C471-bhoga, BN-C471-ogyan-hooya, BN-C471-shojyashayi

\chapter{Falling Ill}
\label{ch:bn-falling-ill}
\hlchaptermodality{471}

\begin{chapteropening}
\textbf{By the end of this chapter:} \emph{I can say that I am ill, that I feel dizzy or that someone fainted.}

Complete the last lesson of chapter 471: I can say that I am ill, that I feel dizzy or that someone fainted.
\end{chapteropening}

\section[\texorpdfstring{\bn{অসুখ} (ôsukh)}{অসুখ (ôsukh)}]{\bn{অসুখ} (ôsukh) — an illness}

\label{lesson:BN-C471-osukh}



\begin{quote}
Before the new one: say the Bengali for to fall out, then the Bengali for a knot of hair.
\end{quote}

\subsection*{You'll want to know: \bn{অসুখ}}
\textbf{\bn{অসুখ}} — \emph{ôsukh} — \textquotedblleft{}an illness\textquotedblright{}.

\begin{quote}
\bn{অসুখের} \bn{জন্য} \bn{অরুণ} \bn{এক} \bn{সপ্তাহ} \bn{স্কুলে} \bn{গেল} \bn{না।}
\end{quote}

\begin{quote}
\emph{ôsukher jônyô ôrun ek sôptāhô skule gelô nā.}
\end{quote}

\begin{quote}
Arun missed a week of school because of illness.
\end{quote}

\bn{অসুখ} \bn{করা}, to fall ill: \bn{আমার} \bn{অসুখ} \bn{করল}.

\begin{cousinweb}
\textbf{\bn{অসুখ}} (\emph{ôsukh}) is Sanskrit \emph{asukha}, \textquotedblleft{}unhappy, painful\textquotedblright{} (Monier-Williams): \emph{a-}, \textquotedblleft{}not\textquotedblright{}, on \emph{sukha}, the \bn{সুখ} (ease) of an earlier lesson. Bengali made it the everyday word for an illness, a dis-ease.
\end{cousinweb}

\begin{grammarlens}[title={What you've built}]
One more word for describing what you lived and felt.
\end{grammarlens}

\subsection*{Your turn}
\begin{itemize}
\raggedright
  \item \emph{Say it:} \emph{ôsukh}
  \item \emph{Say it:} \emph{ôsukh}, once more
  \item \emph{Say it:} \emph{ôsukher jônyô ôrun ek sôptāhô skule gelô nā.} — Arun missed a week of school because of illness.
  \item \emph{Recall:} say the Bengali for to fall out, then the Bengali for a knot of hair, then say \emph{ôsukh} again
\end{itemize}

\begin{tcolorbox}[breakable,colback=teal!4,colframe=teal!35!black,title={Before you move on}]
What does \bn{অসুখ} mean? (\textquotedblleft{}An illness\textquotedblright{}.) Say it once more.
\end{tcolorbox}
\section[\texorpdfstring{\bn{লক্ষণ} (lôkkhôn)}{লক্ষণ (lôkkhôn)}]{\bn{লক্ষণ} (lôkkhôn) — a symptom; a sign}

\label{lesson:BN-C471-lokkhon}



\begin{quote}
Before the new one: say the Bengali for a knot of hair, then the Bengali for an illness.
\end{quote}

\subsection*{You'll want to know: \bn{লক্ষণ}}
\textbf{\bn{লক্ষণ}} — \emph{lôkkhôn} — \textquotedblleft{}a symptom; a sign\textquotedblright{}.

\begin{quote}
\bn{জ্বর} \bn{আর} \bn{কাশি} \bn{এই} \bn{রোগের} \bn{লক্ষণ।}
\end{quote}

\begin{quote}
\emph{jbôr ār kāshi ei roger lôkkhôn.}
\end{quote}

\begin{quote}
Fever and a cough are symptoms of this disease.
\end{quote}

\begin{cousinweb}
\textbf{\bn{লক্ষণ}} (\emph{lôkkhôn}) is a tatsama, Sanskrit \emph{lakṣaṇa}, \textquotedblleft{}a mark, a sign\textquotedblright{}, used also of the signs of a disease (Monier-Williams).
\end{cousinweb}

\begin{grammarlens}[title={What you've built}]
One more word for describing what you lived and felt.
\end{grammarlens}

\subsection*{Your turn}
\begin{itemize}
\raggedright
  \item \emph{Say it:} \emph{lôkkhôn}
  \item \emph{Say it:} \emph{lôkkhôn}, once more
  \item \emph{Say it:} \emph{jbôr ār kāshi ei roger lôkkhôn.} — Fever and a cough are symptoms of this disease.
  \item \emph{Recall:} say the Bengali for a knot of hair, then the Bengali for an illness, then say \emph{lôkkhôn} again
\end{itemize}

\begin{tcolorbox}[breakable,colback=teal!4,colframe=teal!35!black,title={Before you move on}]
What does \bn{লক্ষণ} mean? (\textquotedblleft{}A symptom; a sign\textquotedblright{}.) Say it once more.
\end{tcolorbox}
\section[\texorpdfstring{\bn{ভোগা} (bhogā)}{ভোগা (bhogā)}]{\bn{ভোগা} (bhogā) — to suffer from (an illness)}

\label{lesson:BN-C471-bhoga}



\begin{quote}
Before the new one: say the Bengali for an illness, then the Bengali for a symptom.
\end{quote}

\subsection*{You'll want to know: \bn{ভোগা}}
\textbf{\bn{ভোগা}} — \emph{bhogā} — \textquotedblleft{}to suffer from (an illness)\textquotedblright{}.

\begin{quote}
\bn{দাদু} \bn{তিন} \bn{সপ্তাহ} \bn{ধরে} \bn{জ্বরে} \bn{ভুগছেন।}
\end{quote}

\begin{quote}
\emph{dādu tin sôptāhô dhôre jbôre bhugchhen.}
\end{quote}

\begin{quote}
Grandfather has been suffering from fever for three weeks.
\end{quote}

Past (he/she …): \bn{তিনি} \bn{কাশিতে} \bn{ভুগলেন}, \emph{tini kāshite bhuglen}.

The illness takes -\bn{এ} or -\bn{তে}: \bn{জ্বরে} \bn{ভোগা}, to suffer from fever.

\begin{cousinweb}
\textbf{\bn{ভোগা}} (\emph{bhogā}) is listed by CDIAL for Bengali, with the sense of going through or experiencing something, under a reconstructed form (CDIAL 9523); Bengali uses it above all for suffering an illness.
\end{cousinweb}

\begin{grammarlens}[title={What you've built}]
One more word for describing what you lived and felt.
\end{grammarlens}

\subsection*{Your turn}
\begin{itemize}
\raggedright
  \item \emph{Say it:} \emph{bhogā}
  \item \emph{Say it:} \emph{bhogā}, once more
  \item \emph{Say it:} \emph{dādu tin sôptāhô dhôre jbôre bhugchhen.} — Grandfather has been suffering from fever for three weeks.
  \item \emph{Recall:} say the Bengali for an illness, then the Bengali for a symptom, then say \emph{bhogā} again
\end{itemize}

\begin{tcolorbox}[breakable,colback=teal!4,colframe=teal!35!black,title={Before you move on}]
What does \bn{ভোগা} mean? (\textquotedblleft{}To suffer from (an illness)\textquotedblright{}.) Say it once more.
\end{tcolorbox}
\section[\texorpdfstring{\bn{অজ্ঞান} \bn{হওয়া} (ôgyān hôoyā)}{অজ্ঞান হওয়া (ôgyān hôoyā)}]{\bn{অজ্ঞান} \bn{হওয়া} (ôgyān hôoyā) — to faint; to lose consciousness}

\label{lesson:BN-C471-ogyan-hooya}



\begin{quote}
Before the new one: say the Bengali for a symptom, then the Bengali for to suffer from.
\end{quote}

\subsection*{You'll want to know: \bn{অজ্ঞান} \bn{হওয়া}}
\textbf{\bn{অজ্ঞান} \bn{হওয়া}} — \emph{ôgyān hôoyā} — \textquotedblleft{}to faint; to lose consciousness\textquotedblright{}.

\begin{quote}
\bn{রোদে} \bn{দাঁড়িয়ে} \bn{থেকে} \bn{মহিলাটি} \bn{অজ্ঞান} \bn{হয়ে} \bn{গেলেন।}
\end{quote}

\begin{quote}
\emph{rode dā̃ṛiye theke môhilāṭi ôgyān hôye gelen.}
\end{quote}

\begin{quote}
After standing in the sun, the woman fainted.
\end{quote}

Past (he/she …): \bn{তিনি} \bn{হঠাৎ} \bn{অজ্ঞান} \bn{হলেন}, \emph{tini hôṭhāt ôgyān hôlen}.

\begin{cousinweb}
\textbf{\bn{অজ্ঞান} \bn{হওয়া}} (\emph{ôgyān hôoyā}) is Sanskrit \emph{ajñāna}, \textquotedblleft{}not knowing\textquotedblright{} (Monier-Williams), \emph{a-} on the \bn{জ্ঞান} (taught) of knowledge, with \emph{hôoyā}, \textquotedblleft{}to be, to become\textquotedblright{} (CDIAL 9416): to become senseless.
\end{cousinweb}

\begin{grammarlens}[title={What you've built}]
One more word for describing what you lived and felt.
\end{grammarlens}

\subsection*{Your turn}
\begin{itemize}
\raggedright
  \item \emph{Say it:} \emph{ôgyān hôoyā}
  \item \emph{Say it:} \emph{ôgyān hôoyā}, once more
  \item \emph{Say it:} \emph{rode dā̃ṛiye theke môhilāṭi ôgyān hôye gelen.} — After standing in the sun, the woman fainted.
  \item \emph{Recall:} say the Bengali for a symptom, then the Bengali for to suffer from, then say \emph{ôgyān hôoyā} again
\end{itemize}

\begin{tcolorbox}[breakable,colback=teal!4,colframe=teal!35!black,title={Before you move on}]
What does \bn{অজ্ঞান} \bn{হওয়া} mean? (\textquotedblleft{}To faint; to lose consciousness\textquotedblright{}.) Say it once more.
\end{tcolorbox}
\section[\texorpdfstring{\bn{শয্যাশায়ী} (shôjyāshāyī)}{শয্যাশায়ী (shôjyāshāyī)}]{\bn{শয্যাশায়ী} (shôjyāshāyī) — bedridden; confined to bed (by illness)}

\label{lesson:BN-C471-shojyashayi}



\begin{quote}
Before the new one: say the Bengali for to suffer from, then the Bengali for to faint.
\end{quote}

\subsection*{You'll want to know: \bn{শয্যাশায়ী}}
\textbf{\bn{শয্যাশায়ী}} — \emph{shôjyāshāyī} — \textquotedblleft{}bedridden; confined to bed (by illness)\textquotedblright{}.

\begin{quote}
\bn{অসুখের} \bn{পরে} \bn{দাদু} \bn{এক} \bn{মাস} \bn{শয্যাশায়ী} \bn{ছিলেন।}
\end{quote}

\begin{quote}
\emph{ôsukher pôre dādu ek mās shôjyāshāyī chhilen.}
\end{quote}

\begin{quote}
After his illness my grandfather was confined to bed for a month.
\end{quote}

Another sign of being unwell: \bn{মাথা} \bn{ঘোরা}, to feel dizzy (the head goes round): \bn{আমার} \bn{মাথা} \bn{ঘুরছে}.

\begin{cousinweb}
\textbf{\bn{শয্যাশায়ী}} (\emph{shôjyāshāyī}) puts Sanskrit \emph{śayyā}, \textquotedblleft{}a bed\textquotedblright{}, with \emph{śāyin}, \textquotedblleft{}lying, resting\textquotedblright{} (Monier-Williams): lying in bed, and so laid up by illness.
\end{cousinweb}

\begin{grammarlens}[title={What you've built}]
One more word for describing what you lived and felt.
\end{grammarlens}

\subsection*{Your turn}
\begin{itemize}
\raggedright
  \item \emph{Say it:} \emph{shôjyāshāyī}
  \item \emph{Say it:} \emph{shôjyāshāyī}, once more
  \item \emph{Say it:} \emph{ôsukher pôre dādu ek mās shôjyāshāyī chhilen.} — After his illness my grandfather was confined to bed for a month.
  \item \emph{Recall:} say the Bengali for to suffer from, then the Bengali for to faint, then say \emph{shôjyāshāyī} again
\end{itemize}

\begin{tcolorbox}[breakable,colback=teal!4,colframe=teal!35!black,title={Before you move on}]
What does \bn{শয্যাশায়ী} mean? (\textquotedblleft{}Bedridden; confined to bed (by illness)\textquotedblright{}.) Say it once more.
\end{tcolorbox}
